\chapter{Architecture}
\label{ch:architecture}

The design begins with a constrained composition problem: one exact base, independently installed siblings, and an explicitly selected set. We chose this topology to make each module's parent assumptions checkable and to keep ordinary Ubuntu package paths. The implementation then separates admission, semantic assembly, and execution. Those operations answer different questions, so their success conditions remain distinct.

\section{Pinned sibling model and invariants}

Let B be an exact base artefact and M a selected set of deltas built against it. Each module contributes a package map P(m), removals R(m), captured filesystem changes, and extracted identity metadata. At the package-summary level, the effective installed state is described by

\begin{equation}
P_{effective}(M)=\left(P(B)\cup\bigcup_{m\in M}P(m)\right)\setminus\bigcup_{m\in M}R(m).
\end{equation}
This is a union under compatibility conditions, not permission to overwrite contradictory versions in a dictionary. The admission checker rejects incompatible contributions before treating the expression as one package view. The removal term describes the metadata model; it does not claim that arbitrary file deletions and removal interactions have been validated. The current additive build policy and whiteout guard impose a narrower practical scope.

The generation identifier is derived from \texttt{(snapshot, suite, arch, base\_sha256)}. Module specifications and module versions are excluded so a module may change within that generation. A changed base at the same snapshot is nevertheless a different generation. The identifier provides a checkable compatibility boundary; because it is an unkeyed digest stored with the bundle, it cannot authenticate an origin against deliberate replacement and resealing.

\Cref{fig:architecture} locates the boundaries between building a reusable bundle and constructing a selected machine image. The reconciliation layer is generated for a selected set; it is not another independently built workload module.
\begin{figure}[tbp]
\centering
\begin{tikzpicture}[node distance=9mm and 9mm, box/.style={draw,rounded corners,align=center,text width=3.55cm,minimum height=11mm,font=\small}, >={Stealth}]
\node[box] (base) {Pinned archive\\exact base};
\node[box,right=of base] (deltas) {Independent builds\\sibling bundles};
\node[box,right=of deltas] (admit) {Selected set\\metadata admission};
\node[box,below=of admit] (merge) {Overlay and reconcile\\physical verification};
\node[box,left=of merge] (pack) {Add boot inputs\\flatten and pack};
\node[box,left=of pack] (boot) {Boot image\\observe services};
\draw[->] (base)--(deltas);\draw[->] (deltas)--(admit);\draw[->] (admit)--(merge);\draw[->] (merge)--(pack);\draw[->] (pack)--(boot);
\end{tikzpicture}
\caption{Server-side construction and verification of a selected root. Schematic; no measured quantities are represented. Diagram prepared with tool H1.}
\label{fig:architecture}
\end{figure}

\section{Conflict taxonomy and policy}

The conflict taxonomy separates mechanisms by the action they require. Some overlaps are harmless; some states must be rejected; others should be prevented before the bytes are created; shared registries need reconciliation; runtime resource contention needs execution. The policies in the taxonomy distinguish the observed mechanisms.

\Cref{tab:taxonomy} summarizes the corresponding contract.
\begin{table}[tbp]
\centering\small
\begin{tabularx}{\textwidth}{@{}p{2.6cm}YY@{}}
\toprule
Class & Mechanism & Policy \\
\midrule
\num{1}. Benign overlap & Siblings contribute the same package at the same version. & Accept at the package level and measure duplication. \\
\num{2}. Version skew & Package versions or parent/snapshot/generation identities disagree. & Reject before mounting. \\
\num{3}. Declared conflict & \texttt{Conflicts} or \texttt{Breaks}, including resolved virtual names, forbids coexistence. & Reject through package-relation checking. \\
\num{4}. File collision & Different packages own one path without a sanctioned replacement or diversion. & Intersect ownership sidecars and reject unsanctioned overlap. \\
\num{5}. State divergence & Independent installations rewrite complete shared registries. & Merge records or regenerate derived state. \\
\num{6}. Implicit base upgrade & A delta upgrades a package inherited from base. & Upgrade base first at the pin and detect recurrence. \\
\num{7}. Identity collision & Numeric owners and account identities disagree across siblings. & Partition allocation windows, merge records, and audit owners. \\
\num{8}. Runtime resource conflict & Installed services compete for a runtime resource. & Observe at Tier \num{3}; the earlier tiers do not model actual resource acquisition. \\
\num{9}. Opaque directory erasure & A directory marker captured in one build hides another sibling's files. & Strip opaque markers under checked build preconditions; check visibility. \\
\bottomrule
\end{tabularx}
\caption{Conflict mechanisms and their composition policies. The classes describe mechanisms, not mutually exclusive set outcomes.}
\label{tab:taxonomy}
\end{table}


\subsection{Package compatibility and file ownership}

Class \num{1} treats matching package identity as benign overlap for admission; it does not prove that post-install generated bytes match. Class \num{2} detects incompatible package versions and composability preconditions. Class \num{3} evaluates declared relations, including virtual names. A provided capability is not intrinsically exclusive: two providers without an explicit conflict are not rejected merely for providing the same name.

Module-level requirements are checked separately from Debian package relations. A selected CUDA module can require a selected driver module and an acceptable module version even if each was built alone against base. A set missing its driver may become acceptable when that driver is added. This is why a pairwise rejection is not automatically a permanent exclusion edge for a larger set. The current admission census demonstrates that distinction.

Class \num{4} needs the path-ownership sidecars, because package metadata alone does not enumerate every collision. Ownership intersection followed by replacement and diversion analysis has a direct predecessor in EDOS/Mancoosi \autocite[\S~2.2.5]{treinen2008solving}. \texttt{Replaces} and diversions require special treatment, but a declaration allowing dpkg to overwrite a path during sequential installation does not by itself reproduce those semantics in a static overlay. ModFS uses a conservative admission policy; complete replacement and diversion behavior remains a limit to the generality of the result. A package-owned path index also cannot classify every inode attribute or unowned generated file.

\subsection{Shared state and numeric identities}

Class \num{5} requires a semantic result for each shared database. Merely exempting its path from a collision report stops a warning without preserving the records. The registry-specific operations below define what is combined and what is regenerated.

Class \num{7} also requires prevention. Package scripts in independent roots can allocate the same free number to different services. Once files are stored with that owner, merging names does not remap the inodes. Permanent allocation windows prevent ordinary builds from making the collision. Both allocator interfaces must be configured before installation, and account metadata and file-owner lists must be checked afterwards. Shared accounts supplied by base still require consistent identities. The final deployed allocation policy must also avoid reusing module IDs for future local users, an unresolved lifecycle concern.

\subsection{Context-sensitive filesystem and runtime effects}

Union-mount research already describes opaque-directory semantics and the sensitivity of deletion markers to layer order \autocite{pendry1995union}. Class \num{9} arose because a captured filesystem operation has meaning relative to its lower layers. During a build, an opaque marker can hide only base entries. During composition, the same marker can hide a sibling that was absent when it was created. In the motivating case, \texttt{pyyaml} hid the \texttt{pytools} numpy tree without two packages claiming the same file path. A successful representation round trip through SquashFS therefore did not prove the desired sibling semantics.

The current build policy strips opaque markers only for direct-base siblings and rejects upperdirs containing whiteouts, a conservative guard against a deleting module whose directory-replacement intent could be destroyed. The storage format's ability to preserve whiteouts is not a validated arbitrary-removal workflow. V7 checks the resulting visible paths but still leaves equal-size content substitution outside its predicate.

Class \num{8} was added after boots exposed a failure beyond the successful package and structural checks. It marks the runtime boundary: \texttt{apache2} and \texttt{nginx} can be metadata-compatible and structurally complete while both seek port \num{80}. Their actual startup outcome depends on scheduling. The existing Tier-1 and Tier-2 models cannot observe that race. A future policy could statically warn about known configured listeners, but would not turn an offline package or path check into an observation of resource acquisition. Runtime results are therefore recorded separately from structural success.

\section{Building against the exact parent}

The base builder uses \texttt{debootstrap}, then upgrades the base against the same snapshot pockets that sibling installations will see. This prevents a module from receiving newer base libraries merely because its \ac{APT} view includes updates absent from the initial bootstrap. The snapshot identity is in the archive URL so both bootstrap and subsequent \ac{APT} operations use the same archive view.

A sibling build mounts the exact parent \texttt{base.sqsh} read-only and uses it as the lower filesystem of an OverlayFS mount. Packages are installed into the merged view, with an upperdir capturing newly written files, copies of modified base files, and deletion markers. The builder compresses only that upperdir. The exact-parent requirement matters: a leftover unpacked \texttt{base.dir} could contain bytes different from the archive whose digest is recorded in the generation. The immutable archive therefore supplies the parent; a leftover unpacked build directory is insufficient.

Before \ac{APT} runs, both account allocators are redirected to the assigned window: \texttt{adduser.conf} controls \texttt{adduser --system}, and \texttt{login.defs} controls \texttt{useradd -r}. Both are needed because package scripts use both tools. The builder restores the original allocator configuration before squashing. Deleting a base file instead would leave a whiteout that hides the base policy during composition. The extractor later checks declared accounts and numeric owners of shipped files, because a file can borrow another module's \ac{UID} even when its own module never declared that account.

Build-only mounts and daemon-start suppression let maintainer scripts run in the target root without treating it as a booted machine. A \texttt{policy-rc.d} hook prevents service startup; apt/dpkg logs and transient caches are excluded; the base machine identity is cleared; and archive inode timestamps are pinned. Completion markers are written after successful work. These measures remove specific unwanted outputs, but do not erase a password-change date embedded in \texttt{/etc/shadow}, a generated database identifier, or a host-dependent post-install result. They must not be described as a proof of byte reproducibility.

The builder also removes \texttt{trusted.overlay.opaque} under the checked flat-sibling preconditions. Such a marker was created against the base but would later hide entries from a sibling that did not exist at build time. The storage representation preserves ordinary file whiteouts, but the current builder conservatively refuses an upperdir containing them before applying this opaque-marker policy. Thus representation support for deletion is not demonstrated support for building arbitrary deleting siblings. The evaluated pattern is additive; removals and directory replacement remain a separate validation problem.

\section{Bundle representation and integrity}

\subsection{Artefact, manifest, and sidecar}

Each module consists of \texttt{<name>.sqsh}, \texttt{<name>.json}, and \texttt{<name>.files.json.zst}. The manifest is small enough for routine admission; the sidecar supplies per-path ownership when Class \num{4} needs it; physical composition mounts the archive. Thus Tier \num{1} avoids mounts and root privileges, although it does consume the relevant sidecars as well as manifests. It does not rehash all SquashFS bytes on every metadata decision: artefact consumption and re-derivation have their own checks.

\Cref{tab:manifest} summarizes the corresponding contract.
\begin{table}[tbp]
\centering\small
\begin{tabularx}{\textwidth}{@{}p{2.6cm}Y@{}}
\toprule
Manifest group & Fields and meaning \\
\midrule
Identity & \texttt{schema}, \texttt{module}, \texttt{version}, \texttt{parent}, \texttt{snapshot}, \texttt{suite}, \texttt{arch}, \texttt{built}, \texttt{generation}. \\
Installation & \texttt{requested}, contribution \texttt{packages}, and \texttt{removed}; each package retains its version and raw Debian relations. \\
Module policy & \texttt{requires}, \texttt{conflicts}, \texttt{provides}. \\
Identity and service observations & \texttt{uid\_range}, \texttt{accounts}, \texttt{identity\_audit}, \texttt{units}. \\
Artefact identity & \texttt{artifact.file}, \texttt{artifact.bytes}, \texttt{artifact.sha256}. \\
Integrity policy & \texttt{binding}, including required field coverage, source, and digests. \\
\bottomrule
\end{tabularx}
\caption{Responsibilities of the module manifest. Requested packages describe intent; extracted contributions describe the stored installation.}
\label{tab:manifest}
\end{table}


The sealed field set is explicit; ancillary metadata is not assumed to be covered merely because it appears in the document.

A module records only package contributions that differ from its parent, plus removals. Consumers reconstruct an effective package map from the base and contributions. If two selected sources offer incompatible versions, admission rejects the set before a map overwrite could conceal the disagreement. Package relations and module relations remain distinct: the latter require a selected module name or provided capability at an acceptable module version. Multiple providers of one capability are not automatically rejected unless an explicit conflict expresses that policy; provider cardinality is not represented.

The sidecar maps paths to owning package names and records diversions. Extraction operates on the stored module and parent artefacts so excluded build scratch cannot be mistaken for shipped content. This also lets the build tree be discarded. A sealed manifest can still describe the wrong \emph{requested specification} if no one compares the spec with the installation; \texttt{jq}'s earlier missing \texttt{moreutils} is the project's concrete example. Bundle consistency and fulfilment of operator intent are distinct conditions.

\subsection{What each digest proves}

\Cref{tab:digests} summarizes the corresponding contract.
\begin{table}[tbp]
\centering\small
\begin{tabularx}{\textwidth}{@{}p{2.6cm}YY@{}}
\toprule
Commitment & Bytes or fields covered & Consumer obligation \\
\midrule
Artefact digest & The exact \texttt{.sqsh} bytes, named by \texttt{artifact.sha256} and repeated in \texttt{binding.artifact\_sha256}. & Check agreement of the recorded fields and hash the actual archive when consuming its bytes. \\
Sidecar digest & \textbf{Decompressed \ac{JSON} bytes} from \texttt{.files.json.zst}. & Decompress, verify \texttt{binding.sidecar\_sha256}, then parse and use ownership data. \\
Field digest & Canonical \ac{JSON} of the required, named manifest fields. & Require the complete field list and presence of its members, then recompute \texttt{binding.fields\_sha256}. \\
\bottomrule
\end{tabularx}
\caption{Integrity contracts for a filesystem bundle. Each consumer must require the appropriate commitment before using its content.}
\label{tab:digests}
\end{table}


Changing the zstd encoding without changing its decompressed \ac{JSON} is outside the sidecar digest's scope.

The field digest includes the \texttt{artifact} record and \texttt{generation}. The field list is required by the consumer, not negotiated by the document being checked. Otherwise a shortened document could announce a shorter list and pass its own hash. The sidecar digest is required separately because the \texttt{binding} object is not part of the field payload. This is a schema-policy issue, not an inherent cryptographic rule that a sidecar digest could never itself be covered by another digest.

The binding verifier mounts artefacts and re-derives metadata, which is stronger than checking that a document's hashes agree internally. Even re-derivation does not establish the identity of a trusted publisher. Someone able to replace the whole bundle and reseal it can create a consistent different bundle. Signing and an external trust root are outside scope.

A generation consists of snapshot, suite, architecture, and exact base digest, together with an identifier derived from those values. Module specifications are deliberately absent: a specification may change within one generation. A metadata refresh preserves the recorded generation rather than relabelling old bytes as children of the current base. Adoption for legacy artefacts is an explicit assertion about their origin, not reconstruction of missing historical proof.

\section{Registry reconciliation}

The composer mounts each selected SquashFS layer and overlays a writable reconciliation layer above them. It supplies the reconciler with layer roots in the intended precedence order. The reconciler writes parsed registry results into the top layer, then the native tools regenerate derived state. Relative lowerdir paths keep mount-option length manageable at high N. Mount creation, reconciliation, regeneration, and teardown must all preserve an observable failure rather than allowing partial scratch state to become a publishable result.

\subsection{Package status and installation marks}

\texttt{dpkg/status} uses blank-line-separated stanzas. The format-level identity is package plus architecture, and different versions of the same effective package are incompatible. The current implementation, in the documented single-architecture scope, indexes status by package name. Consequently it should not be presented as a general multiarch union; an architecture-aware key is a required extension for that claim. Repeated installed stanzas at the same version retain a selected source stanza, while version disagreement contributes an error. Fields outside the compared version can still depend on the selected layer.

\texttt{extended\_states} also uses stanzas, but its semantics differ. \ac{APT} generally records automatic installation positively; manual installation can be represented by the \textbf{absence} of an automatic stanza. Combining only explicit entries would miss a sibling's manual request. The reconciler considers the installed package records together with automatic marks, lets manual intent dominate, and inherits base marks when a delta has no replacement file. An empty output can therefore be necessary to clear a top layer's obsolete automatic flags. This is not ordinary last-layer precedence.

\subsection{Alternatives and diversions}

An alternatives registry is positional: mode, generic link, slave-link definitions, then candidate paths, priorities, and slave targets. It must be parsed with that structure. \texttt{vim} offers \path{/usr/bin/vim.basic} at priority \num{30} and \texttt{emacs} offers \texttt{/usr/bin/emacs} at priority \num{0}; retaining both candidates and running automatic selection chooses vim. If layers assign different priorities to the \emph{same} candidate, the merger reports a conflict. It may retain the higher priority in the scratch output, but its nonzero result means that tree is not accepted. A diagnostic choice is not successful conflict resolution.

A diversion is a consecutive triple: original path, diverted path, and responsible package. Malformed or partial triples must fail rather than silently disappearing. A diversion is identified semantically by its original path. The present diversion parser deduplicates whole triples; it does not independently reject different triples sharing the same original path. Malformed-record rejection does not establish complete conflicting-diversion arbitration. The reference key and implementation behavior therefore differ, limiting the claimed scope.

\subsection{Accounts and debconf}

The account group comprises \texttt{passwd}, \texttt{group}, \texttt{shadow}, \texttt{gshadow}, \texttt{subuid}, and \texttt{subgid}. Their record shapes are centralized in a shared account-format schema; names identify the ordinary account records, while subordinate ranges are retained as whole-line sets. Numeric \ac{UID}, \ac{GID}, and primary-group disagreement is an error. Group and gshadow membership is unioned as sets rather than overwritten. Other fields follow the declared highest-layer policy, and sensitive modes and ownership are preserved. Disjoint allocation windows prevent ordinary numeric collisions, but do not remove the need for this record union.

Debconf's \texttt{config.dat}, \texttt{templates.dat}, and \texttt{passwords.dat} contain records keyed by \texttt{Name}. \texttt{Owners} is a set, so overlapping owners are deduplicated; continuation fields in templates must survive parsing. Conflicting values or other incompatible shared fields are reported, and \texttt{passwords.dat} retains its sensitive attributes. Missing record keys are errors. The parser cannot legitimately call a merge successful merely because it recognized a subset of the source.

The reconciler accumulates problems and returns failure even if it has already written diagnostic output to the scratch tree. Its caller must discard a failed composition. Tested order reversals preserve semantic record sets for the documented examples, but record order and bytes differ, and conflicting highest-layer fields do not acquire general order independence from this mechanism.

\subsection{Derived state and verification coverage}

\texttt{/etc/alternatives/*} is regenerated by \texttt{update-alternatives --auto} from the merged candidates. \texttt{/etc/ld.so.cache} is regenerated by \texttt{ldconfig} from the composed library tree. These are the last two of the eight registry groups. They are functions of the merged inputs, so unioning their old outputs is inappropriate. Failure of either native tool is a failed composition.

V2 checks package status, V3 alternatives, V4 the linker result, V6 accounts, and V8 debconf. V5 asks dpkg to audit its view, while V7 checks path visibility and selected file properties. This coverage does not give every registry an equally strong independent oracle: diversions and extended states do not have dedicated V-numbered semantic checks in the generated list. Trigger registrations also remain outside the handled registry set. These limits matter when interpreting a clean V1--V8 row.

\section{How the registry inventory was established}

The inventory grew through compositions whose visible behavior disagreed with their apparent package state. Lost package records motivated status reconciliation. A survey of paths shared by independent upperdirs exposed alternatives, diversions, and automatic-installation marks; inspecting library lookup exposed the derived cache. Account auditing revealed that unique numeric allocations still needed a union of account records. Enumerating shared unowned files then exposed configuration databases whose answers had disappeared beneath a later layer.

This supplies a repeatable discovery method: intersect the files present in multiple upperdirs and remove package-owned paths using the ownership inventory. The remaining shared files are candidates for semantic interpretation. They are not automatically interchangeable registries, and the survey does not prove complete coverage of every maintainer-script effect. Trigger registrations, for example, remain outside the implemented merge groups. New package families should prompt another survey, not an assumption that the existing list is universal.

\section{From admitted sets to a running image}

A requested plan maps layer counts to sample counts. The sampler combines declared requirements with exclusions learned from Tier-1 observations. It first explains a pair rejection that a later provider could satisfy; only irreducible exclusions belong in the pair-conflict model. Otherwise \texttt{fake-cuda} would be systematically excluded from high-layer sets despite belonging to an admissible set with its driver.

When the admissible space or its complement can be enumerated, \texttt{exact} mode samples uniformly from it and reports its cardinality. For larger spaces, \texttt{sampled} mode uses a seeded, closure-aware greedy construction and is reproducible without being uniform. Every proposal still goes through the admission checker. Refused proposals are evidence of model/checker disagreement, not observations to drop. The plan, seed, mode, attempted sets, and module coverage are needed to interpret a sweep; aggregate timing means alone cannot recover them.

Boot construction begins with bundle validation and admission. An explicit known-negative run may bypass the ordinary admission requirement for a stated experiment, but must not be labelled admitted. The same reconciliation mechanism then assembles the selected root. The packing stage adds a bootable kernel image, creates \ac{GPT}/\ac{EFI} and ext4 filesystems, copies the composed tree with numeric identities and attributes, and writes the boot configuration. Siblings do not contain the boot kernel image, although the driver sibling does contain kernel module objects; calling every sibling a purely userspace delta would conceal that \ac{ABI} dependency.

The guest harness records systemd state, foreign pending jobs, failed units, package audit, per-module probes, and listeners through serial markers. Its wait must exclude its own job, or the observer prevents the steady state it is waiting for. Probe observations prefixed for the harness are retained even on success. Partition devices must be usable before filesystem creation, and the disk must be sized for the actual root contents. These details determine whether the experiment tests the composed system or only a failure in its apparatus.

The bundle retains the serial log, verdict, run identity, and resolved kernel information. Recording a kernel digest establishes which boot input was used; it does not by itself make the whole build reproducible. Old bundles remain immutable, while later evidence generation assesses their currency against the artefact inventory. Freshness checks use recorded identities and rebuild times. A freshness label does not constitute new content verification of every old run.

\section{Verification and trust boundaries}

The checker must declare what it has observed. Missing sidecars or malformed nested ownership fields cannot be treated as an empty, clean set of paths or owners. Where the checker returns \texttt{ACCEPT WITH WARNINGS}, that is weaker than fully checked acceptance. A negative control is likewise specific: an old-snapshot control exercises composability rejection, but does not establish that every account or debconf predicate is effective. Adversarial tests target the individual observation boundaries.

Trusted inputs are another design assumption. Packing and reconciliation invoke tools from composed roots with host-side privileges before the VM boundary. Bundle digests detect drift and internal inconsistency; they do not isolate hostile package code. The prototype therefore concerns a controlled build catalogue, not a service safely composing arbitrary untrusted customer artefacts. This limit follows from where code executes, independently of whether its files have consistent hashes.

\section{Updates, publication, and rollback}

An intra-generation specification update rebuilds the affected sibling against the unchanged exact base, refreshes its manifest and sidecar, and repeats binding, admission, and relevant later checks. Rebuilding at a fixed snapshot cannot obtain a security version absent from that archive. A pin move instead builds a new base and all siblings in an isolated generation. The selected generation must be exposed coherently; the current implementation has no transactional remote publisher providing that atomicity.

Build work and transfer opportunity are separate. A byte-identical artefact could be reused from a content-addressed cache, while a changed hash requires the whole new archive under the current reporting model. An unchanged file size is not sufficient for reuse. No implemented remote client or block-difference protocol turns these byte comparisons into a measured network result.

The selected root is flattened into an ordinary image for delivery. Rollback means redeploying a retained complete bootable image and its configuration identity, or retaining every input needed to reconstruct that image. Base and deltas alone omit the packing-time kernel, initramfs, bootloader, and configuration. Operating-system rollback also does not undo application data or schema changes. These boundaries determine what \cref{ch:analysis} measures and why repository savings cannot be presented as per-node update or rollback performance.
